\section{Introduction}

Rice is one of the most important staple crops worldwide and plays a central role in food security, particularly in regions where it represents a major source of daily caloric intake. The Food and Agriculture Organization reports that rice is the staple food for more than half of the world's population, and that in Asia more than two billion people obtain a large share of their caloric intake from rice and rice-derived products \cite{fao_rice_conference}. Therefore, the early identification of rice diseases is relevant not only for agronomic management but also for protecting crop productivity in rice-dependent regions.

Rice production is affected by several foliar diseases, including blast, bacterial blight, brown spot, and tungro. The International Rice Research Institute identifies disease damage as a factor that can greatly reduce rice yield and describes tungro as one of the most destructive diseases of rice in South and Southeast Asia \cite{irri_rice_diseases,irri_rice_disease_damage}. These diseases manifest through visual symptoms such as spots, lesions, discoloration, necrosis, or leaf deformation, making image-based diagnosis a feasible computer vision task. However, visual symptoms may also vary with plant stage, lighting, image acquisition conditions, and disease severity, which can complicate automatic classification.

Traditional plant disease diagnosis commonly depends on manual inspection, field observation, or expert assessment. Although this process is useful in practice, it can be time-consuming, labor-intensive, subjective, and difficult to scale for large agricultural areas. As a result, computer vision and machine learning methods have become increasingly relevant for supporting automatic crop disease recognition. Public image datasets, such as the Rice Leaf Disease Images dataset used in this study, provide labeled samples for bacterial blight, blast, brown spot, and tungro, enabling supervised experimentation with deep learning models \cite{kaggle_rice_leaf_disease}.

Convolutional neural networks and transfer learning have shown strong performance in plant disease recognition. In rice disease classification, architectures based on DenseNet121 and other deep models have reported high predictive performance \cite{ismail2026densenet}. Lightweight detection models have also been proposed to balance performance and computational cost, as in RiceDetect-Net, which focuses on real-time rice disease detection \cite{yuan2026ricedetectnet}. These studies demonstrate that deep visual representations can capture discriminative patterns from diseased leaves. Nevertheless, high in-distribution accuracy does not always imply that a model is using only disease-related visual cues. Deep neural networks may exploit shortcut features or spurious correlations that perform well on a benchmark but generalize poorly under more challenging conditions \cite{geirhos2020shortcut}. In plant disease datasets, this concern is especially relevant because background, acquisition conditions, and dataset-specific artifacts may correlate with class labels \cite{noyan2022plantvillage_bias}.

Several research lines attempt to reduce these limitations by incorporating segmentation or explicit handcrafted descriptors. Segmentation-based approaches aim to isolate the leaf or lesion region so that the classifier focuses on more relevant visual information. For example, contour-driven segmentation has been used to improve rice leaf disease classification with transfer learning architectures \cite{shakeel2026contour}, while automatic segmentation and multi-scale feature fusion have been explored for plant disease detection in complex backgrounds \cite{arora2025complex}. In a different but complementary direction, explicit texture descriptors such as gray-level co-occurrence matrix (GLCM), local binary patterns (LBP), color histograms, and frequency-domain descriptors have been used to represent local texture, spatial intensity relationships, and color distributions. Feature fusion using GLCM, LBP, and color histograms has been proposed for rice leaf disease classification \cite{sutikno2026glcm_lbp_color}, supporting the idea that handcrafted descriptors can still provide useful information for disease recognition.

However, the reviewed literature suggests that these components are often studied separately, embedded in complex architectures, or evaluated without isolating their individual and combined contributions under the same lightweight CNN protocol. CNN-based works tend to prioritize final predictive accuracy, segmentation-based works often focus on region extraction or complex pipelines, and texture-based works frequently rely on handcrafted feature fusion or traditional classifiers. Consequently, there remains a practical methodological gap: it is not sufficiently clear whether classical segmentation, explicit texture descriptors, or their combination with CNN embeddings provide measurable benefits over a lightweight CNN baseline when evaluated under a controlled ablation study for rice leaf disease classification.

To address this gap, this paper evaluates four configurations using a common dataset, split strategy, training setup, and evaluation protocol: (i) a MobileNetV2 baseline trained on original RGB images, (ii) a CNN trained on segmented images, (iii) a CNN fused with explicit texture features, and (iv) a hybrid configuration combining segmentation and texture features.

The main contributions of this work are summarized as follows:
\begin{itemize}
    \item A controlled ablation study comparing RGB input, segmentation, explicit texture features, and segmentation-texture fusion in lightweight CNN-based rice leaf disease classification.
    \item A leakage-aware dataset partitioning procedure based on perceptual-hash grouping to avoid visually equivalent samples across training, validation, and test sets.
    \item A multi-seed evaluation protocol reporting mean and standard deviation for macro-F1, accuracy, precision, recall, inference time, parameter count, and model size.
    \item A critical interpretation of whether segmentation and texture features improve predictive performance or mainly add complexity under a high-performing lightweight CNN baseline.
\end{itemize}

The remainder of this paper is organized as follows. Section II reviews related work on rice leaf disease classification, segmentation-based recognition, and explicit texture descriptors. Section III describes the research design, dataset, leakage control, experimental scenarios, and evaluation protocol. Section IV presents the preprocessing, segmentation, feature extraction, model architecture, and implementation details. Section V reports the experimental results. Section VI discusses the findings, limitations, and methodological implications. Finally, Section VII presents the conclusions and future work.

